
%(BEGIN_QUESTION)
% Copyright 2013, Tony R. Kuphaldt, released under the Creative Commons Attribution License (v 1.0)
% This means you may do almost anything with this work of mine, so long as you give me proper credit

Et av de unike trekkene ved dette programmet er bruk av {\it mastery-prøver}, der studentene må svare 100\% riktig for å bestå. Vanlige ``proporsjonale'' prøver lar elever bestå ved å oppnå en minimumsscore. Problemet med denne strategien er at elever kan bestå uten å ha lært {\it alle} begrepene de skal kunne. Hensikten med mastery-prøver er å sikre mestring av {\it alle} kritiske begreper ved å kreve 100\% riktighet.

\vskip 10pt

Læreren deler ut eksemplarer av mastery-prøven for emnet INST200 Introduction to Instrumentation, som dekker viktige begreper i kretsanalyse fra første studieår. Gjør ditt beste for å svare riktig på alle oppgaver. Hvis noe er feil på første forsøk, markerer læreren hvilke {\it deler} (ikke hvilke {\it oppgaver}) som er feil, og du får ett nytt forsøk. Hvis den ikke bestås på andre forsøk, regnes det som stryk.

\vskip 10pt

Mastery-prøver kan tas på nytt flere ganger uten karaktertrekk. Hensikten er å gi elevene konstruktiv tilbakemelding og trening slik at alle begrepene mestres. Likevel må prøven bestås innen en fastsatt frist (vanligvis datoen for neste emneprøve) for å få bestått i emnet.

\vskip 10pt

I INST230 brukes denne mastery-prøven kun som trening. I INST200 (høst-, vinter- og vårkvartal) må den samme prøven bestås i sin helhet for å få bestått i INST200, med flere muligheter for nytt forsøk gjennom kvartalet.

\vskip 10pt

I sommerserien (INST230, 231 og 232) er rele-ladderlogikk svært viktig, og du bør derfor følge spesielt med på oppgave \#9 i INST200 mastery-prøven. Hvis du ikke har mestret hvordan relekretser fungerer og analyseres (samme stoff som i ELTR140 Digital 1 i førsteåret), vil du få betydelige utfordringer med motorstyring og PLC-programmering dette kvartalet.

\vskip 10pt

\underbar{file i01235}
%(END_QUESTION)





%(BEGIN_ANSWER)

 
%(END_ANSWER)





%(BEGIN_NOTES)


%INDEX% Course organization, assessment: mastery (practice)
%INDEX% Mastery exam practice

%(END_NOTES)

